\documentclass[11pt]{amsart}
\usepackage[T1]{fontenc}
\usepackage[utf8]{inputenc}
\usepackage[french]{babel}
\usepackage{amssymb}
\usepackage[margin=1.05in]{geometry}
\usepackage[hidelinks]{hyperref}
\newcommand{\stamptime}{\the\numexpr\time/60\relax:\ifnum\numexpr\time-60*(\time/60)\relax<10 0\fi\the\numexpr\time-60*(\time/60)\relax}
\AddToHook{shipout/foreground}{\put(0.5\paperwidth,-\paperheight+0.55cm){\makebox[0pt]{\tiny\ttfamily unit/doc/\jobname.tex \ --- \ compiled \today, \stamptime}}}
\newtheorem{theorem}{Théorème}
\newtheorem{proposition}[theorem]{Proposition}
\newtheorem{lemma}[theorem]{Lemme}
\newtheorem{corollary}[theorem]{Corollaire}
\theoremstyle{definition}
\newtheorem{definition}[theorem]{Définition}
\theoremstyle{remark}
\newtheorem{remark}[theorem]{Remarque}
\newcommand{\pgcd}{\operatorname{pgcd}}
\newcommand{\cont}{\mathrm c}
\title{Anneaux à pgcd : $A$ à pgcd $\Rightarrow$ $A[X]$ à pgcd}
\date{26 octobre 2026}
\begin{document}
\maketitle

\noindent\textbf{Statut.} Preuve informelle écrite avant la
formalisation de \texttt{sag/gcd} (Mines--Richman--Ruitenburg, ch.~IV,
lemme de Gauss par le contenu), dans la forme qui est formalisée. But :
la donnée (S0)/(S1) de \texttt{acl-zmt-formalisation.tex}
(section~\ref{sec:S}). L'état de la formalisation est dans
\texttt{sag/gcd/README.txt}.

\section{Définitions}

Tous les anneaux sont commutatifs et sont des ensembles.

\begin{definition}
\begin{enumerate}
\item $A$ est \emph{intègre discret} si l'on sait décider $x=0$ pour
  tout $x$, si $1\ne0$, et si $x\ne0$, $y\ne0$ entraînent $xy\ne0$.
  ($1\ne0$ sert au lemme~\ref{lem:elim}, où $\lambda=1$ doit être non
  nul.)
\item $a\mid b$ est la \emph{donnée} d'un $q$ avec $b=aq$. On ne
  suppose pas la divisibilité décidable, et elle n'est jamais testée.
\item Un \emph{pgcd} de $a,b$ est la donnée de $d$ avec $d\mid a$,
  $d\mid b$ et, pour tout $c$, $c\mid a$ et $c\mid b$ donnent $c\mid d$.
  Un \emph{anneau à pgcd} est un anneau intègre discret muni d'une
  fonction qui donne un pgcd de tout couple.
\item $f\in A[X]$ est \emph{primitif} si tout diviseur commun de ses
  coefficients divise $1$.
\end{enumerate}
\end{definition}

Un pgcd n'est unique qu'à un inversible près : c'est une donnée, pas une
proposition. On a $\pgcd(0,0)=0$ ($0\mid0$ donne $0\mid d$). Un corps
discret ($x=0$ ou $x$ inversible, et $1\ne0$) est un anneau à pgcd :
$\pgcd(a,b)=b$ si $a=0$, $a$ sinon.

\section{Le contenu}

\begin{lemma}[multiplicativité]\label{lem:M}
Si $d\mid ua$ et $d\mid ub$, alors $d\mid u\,\pgcd(a,b)$.
\end{lemma}
\begin{proof}
Si $u=0$, $u\,\pgcd(a,b)=0$. Sinon soit $g=\pgcd(ua,ub)$ ; $d\mid g$.
Comme $u\mid ua$ et $u\mid ub$, $g=ut$. De $ut\mid ua$ on tire $t\mid a$
(simplification par $u\ne0$), de même $t\mid b$, donc
$t\mid\pgcd(a,b)$ et $g=ut\mid u\,\pgcd(a,b)$.
\end{proof}

Pour une suite $e_0,e_1,\dots$ de $A$ on pose $\cont_0(e)=0$,
$\cont_{n+1}(e)=\pgcd(e_n,\cont_n(e))$.

\begin{lemma}\label{lem:cont}
$\cont_n(e)\mid e_i$ pour $i<n$ ; un diviseur commun des $e_i$, $i<n$,
divise $\cont_n(e)$ ; si $d\mid ue_i$ pour tout $i<n$, alors
$d\mid u\,\cont_n(e)$.
\end{lemma}
\begin{proof}
Récurrence sur $n$ ; la troisième assertion par le lemme~\ref{lem:M}.
\end{proof}

Pour $f\in A[X]$ de degré $<n$ (coefficients nuls à partir de $n$),
$\cont(f)=\cont_n(f)$ divise tous les coefficients, et tout diviseur
commun des coefficients divise $\cont(f)$.

\begin{corollary}\label{cor:P}
Si $f$ est primitif et $d\mid uf_i$ pour tout $i$, alors $d\mid u$.
\end{corollary}
\begin{proof}
$d\mid u\cont(f)$ et $\cont(f)\mid1$.
\end{proof}

\begin{lemma}[partie primitive]\label{lem:pp}
Si $f\ne0$, alors $\cont(f)\ne0$, $f=\cont(f)f^*$ et $f^*$ est primitif.
\end{lemma}
\begin{proof}
Si $\cont(f)=0$, tous les coefficients sont nuls. Les quotients
$f_i=\cont(f)f^*_i$ ($i<n$, et $0$ au-delà) donnent $f^*$. Si $d$ divise
les $f^*_i$, $d\cont(f)$ divise les $f_i$, donc $d\cont(f)\mid\cont(f)$
et $d\mid1$ en simplifiant par $\cont(f)\ne0$.
\end{proof}

\section{Lemme de Gauss}

\begin{lemma}[Gauss]\label{lem:gauss}
Si $E$ est primitif et si $d$ divise tous les coefficients de $Eq$, alors
$d$ divise tous les coefficients de $q$.
\end{lemma}
\begin{proof}
Soit $q=\sum_{j\le m}q_jX^j$. Le théorème de Dedekind--Mertens sous sa
forme de sous-groupes (\texttt{sag/mertens}, sans hypothèse sur $A$) :
pour tout mot $w=(i_0,\dots,i_m)$, le produit $E_{i_0}\cdots E_{i_m}q_j$
est une somme signée de produits $E_{k_1}\cdots E_{k_m}(Eq)_k$. Comme $d$
divise les $(Eq)_k$, il divise donc $E_wq_j$ pour tout mot $w$ de
longueur $m+1$. Par récurrence sur la longueur $n$ : si $d\mid E_wv$ pour
tout mot de longueur $n+1$, alors pour tout mot $w'$ de longueur $n$,
$d\mid E_i(E_{w'}v)$ pour tout $i$, donc $d\mid E_{w'}v$
(corollaire~\ref{cor:P}). À la longueur $0$, $d\mid v$.
\end{proof}

La preuve classique (par un élément premier) n'est pas constructive ; on
n'utilise ni éléments premiers ni factorisation.

\begin{lemma}[division par un primitif]\label{lem:G}
Si $E$ est primitif, $\lambda\ne0$ et $\lambda f=Eq$, alors $E\mid f$.
\end{lemma}
\begin{proof}
$\lambda$ divise les coefficients de $Eq$, donc ceux de $q$
(lemme~\ref{lem:gauss}) : $q=\lambda q'$. Alors $\lambda(f-Eq')=0$ et
$f=Eq'$, coefficient par coefficient.
\end{proof}

\section{L'élimination}

On n'utilise pas le corps des fractions : la division euclidienne dans
$K[X]$ est remplacée par des pas d'élimination à coefficient dominant
près.

\begin{lemma}\label{lem:elim}
Pour $f,g\in A[X]$ on construit $D$ et $\lambda\ne0$ dans $A$ avec
$D\mid\lambda f$, $D\mid\lambda g$, et tels que tout diviseur commun de
$f$ et $g$ divise $D$.
\end{lemma}
\begin{proof}
Récurrence sur $n_f+n_g$, où $f$ (resp.\ $g$) est de degré $<n_f$
(resp.\ $<n_g$). Si $n_f=0$, $f=0$ : $D=g$, $\lambda=1$ ; de même si
$n_g=0$. Sinon, si le coefficient $f_{n_f-1}$ est nul, $f$ est de degré
$<n_f-1$ et l'on récurre ; de même pour $g$. Sinon $\alpha=f_{n_f-1}$ et
$\beta=g_{n_g-1}$ sont non nuls ; si $n_f\ge n_g$, soit $k=n_f-n_g$ et
\[
 f'=\beta f-\alpha X^kg ,
\]
de degré $<n_f-1$ (le coefficient d'indice $n_f-1$ est
$\beta\alpha-\alpha\beta$). La récurrence donne $D,\lambda'$ pour $f',g$.
On pose $\lambda=\lambda'\beta\ne0$ : $\lambda f=\lambda'f'+\alpha X^k
\lambda'g$ et $\lambda g=\beta\lambda'g$ sont divisibles par $D$ ; un
diviseur commun de $f,g$ divise $f'$, donc $D$. Si $n_f<n_g$, on échange
les rôles.
\end{proof}

\begin{theorem}\label{thm:main}
Si $A$ est un anneau à pgcd, $A[X]$ est un anneau à pgcd.
\end{theorem}
\begin{proof}
\emph{Intègre discret.} $f=0$ est décidable (coefficients en nombre fini).
Si $f\ne0$, $g\ne0$ et $fg=0$ : soit $f_i\ne0$ ; Dedekind--Mertens
(conséquence \texttt{ZeroProduct}) donne $f_i^{m+1}g_j=0$, donc $g_j=0$
pour tout $j$, contradiction.

\emph{Pgcd.} Si $f=0$, $\pgcd(f,g)=g$ ; si $g=0$, $f$. Sinon, soit
$D,\lambda$ du lemme~\ref{lem:elim} ; $D\ne0$ (car $\lambda f\ne0$), donc
$D=\cont(D)E$ avec $E$ primitif (lemme~\ref{lem:pp}). Par le
lemme~\ref{lem:G}, $E\mid f$ et $E\mid g$. Soit
$\gamma=\pgcd(\cont(f),\cont(g))$ ; on pose $G=\gamma E$.

$G\mid f$ : $f=Eq$, et $\gamma$ divise les coefficients de $f$, donc
ceux de $q$ (lemme~\ref{lem:gauss}) : $q=\gamma q'$ et $f=Gq'$ ; de même
pour $g$.

Si $h\mid f$ et $h\mid g$ : $h\ne0$, $h=\cont(h)h^*$. Les coefficients de
$f=hr$ sont divisibles par $\cont(h)$, donc $\cont(h)\mid\cont(f)$ ; de
même $\cont(h)\mid\cont(g)$, d'où $\cont(h)\mid\gamma$. D'autre part
$h\mid D=\cont(D)E$, donc $h^*\mid\cont(D)E$ et $h^*\mid E$
(lemme~\ref{lem:G}, $\cont(D)\ne0$). Donc $h=\cont(h)h^*\mid\gamma E=G$.
\end{proof}

\begin{corollary}\label{cor:iter}
Si $k$ est un corps discret, $k[b][a][X][Y][Z]$ est un anneau à pgcd.
\end{corollary}

\section{Les données (S0) et (S1)}\label{sec:S}

\begin{lemma}[Euclide]\label{lem:euclide}
Dans un anneau à pgcd, si $\pgcd(a,b)\mid1$ et $a\mid bc$, alors
$a\mid c$.
\end{lemma}
\begin{proof}
$a\mid ca$ et $a\mid cb$, donc $a\mid c\,\pgcd(a,b)\mid c$
(lemme~\ref{lem:M}).
\end{proof}

\begin{proposition}\label{prop:S}
Soit $S$ un anneau à pgcd, $P,Q\in S$ non tous deux nuls,
$D_0=\pgcd(P,Q)$, $P=D_0P'$, $Q=D_0Q'$. Alors $D_0\ne0$ et
$\pgcd(P',Q')\mid1$ ;
\begin{itemize}
\item[(S0)] $D_0f=0$ entraîne $f=0$ ;
\item[(S1)] si $P'\alpha+Q'\beta=0$ et $P'\ne0$, il existe un unique $w$
  avec $\alpha=-Q'w$, $\beta=P'w$.
\end{itemize}
\end{proposition}
\begin{proof}
Un diviseur commun $e$ de $P',Q'$ donne $D_0e\mid P,Q$, donc
$D_0e\mid D_0$ et $e\mid1$. (S0) : intégrité. (S1) : $P'\mid Q'\beta$,
donc $\beta=P'w$ (lemme~\ref{lem:euclide}) ; puis $P'(\alpha+Q'w)=0$ et
$\alpha=-Q'w$ ; l'unicité par intégrité.
\end{proof}

\section{Points délicats}

\begin{enumerate}
\item \emph{Le pgcd est une donnée}, unique à un inversible près ; une
  fonction pgcd sur l'anneau quotient $A[X]$ (\texttt{PolynomialRing})
  ne se définit pas sur les listes représentantes par descente. On
  utilise le \emph{degré} (le plus petit majorant strict des indices des
  coefficients non nuls), qui est une proposition, donc s'obtient sur le
  quotient ; les coefficients d'indice $<$ degré donnent alors un
  représentant canonique.
\item \emph{Pas d'éléments premiers.} Gauss passe par Dedekind--Mertens et
  le lemme~\ref{lem:M}, qui demande $\pgcd(ua,ub)$ et la simplification
  par $u\ne0$.
\item \emph{Pas de corps des fractions} : élimination avec le facteur
  $\lambda$ (lemme~\ref{lem:elim}).
\item \emph{Divisibilité non décidable} : seul le test $x=0$ est utilisé
  (degré, coefficient dominant, cas $f=0$, $\cont(D)\ne0$).
\item \emph{Corps discret} : la donnée \og$x=0$ ou $x$ inversible\fg{}
  (\texttt{Field.seg}) ne donne l'égalité décidable qu'avec $1\ne0$ ;
  l'hypothèse $1\ne0$ est ajoutée (elle est aussi dans la définition
  d'anneau intègre).
\item \emph{Intégrité de $A[X]$} : elle ne découle pas du lemme de Gauss
  (dans un corps tout élément non nul divise $1$) ; on la tire de
  Dedekind--Mertens par un coefficient non nul $f_i$ et le mot
  $(i,\dots,i)$.
\item \emph{Vers $\Lambda_{XYZ}$} (formalisé, \texttt{sag/gcd}, section suivante).
  La donnée de \texttt{ACLCech.S1T} est dans $\Lambda_{XYZ}$ en
  représentation \texttt{Laurent3} (fonctions $\mathbf Z^3\to A$ à support
  fini), avec $\alpha,\beta,w$ homogènes ; il faut un diviseur homogène
  $D_0$.
\end{enumerate}


\section{Le passage à $\Lambda_{XYZ}$}

Soit $S=A[X][Y][Z]$ et $\psi:S\to\Lambda$ qui envoie $s$ sur la fonction
coefficient ($0$ hors de $\mathbf N^3$).
\begin{enumerate}
\item $\psi$ est additive, $A$-linéaire, envoie la multiplication par
  $X,Y,Z$ sur les décalages, $1$ sur $\delta$, et est injective
  (coefficients, trois fois).
\item Tout $s$ a une liste canonique de termes (coefficients sous le
  degré, aux trois niveaux) $\ell(s)$ avec
  $\psi(ss')=\ell(s)\cdot\psi(s')$.
\item Si $f$ est à support fini (liste $g$) et si $K$ majore les parties
  négatives des exposants de $g$, $X^KY^KZ^Kf=\psi(s)$ pour un $s\in S$.
\item \emph{Régularité.} Si $s\ne0$ et $\ell(s)f=0$ ($f$ à support fini),
  alors $f=0$ : décaler $f$ dans $S$, $S$ est intègre, $\psi$ injective.
  D'où (S0), et (S1) : décaler $\alpha,\beta$ par le même $K$, appliquer la
  proposition~\ref{prop:S} dans $S$, redécaler $w$ ; $w$ est unique, donc
  obtenu sur les troncations du support fini.
\item \emph{Homogénéité.} Composantes $s_n\in S$ ($\psi(s_n)$ est
  $\psi(s)$ en degré $n$) ; le test $s_n=0$ est décidable dans $S$ ; les
  degrés extrêmes $T$ et $B$ de $s\ne0$ par recherche. Pour $x=h+l$,
  $y=h'+l'$ ($h,h'$ composantes de degré maximal),
  $xy=hh'+hl'+ly$, et en degré $T+T'$ seul $hh'\ne0$ contribue ; de même en
  bas. Si $xy$ est homogène de degré $d$, $T_x+T_y=d=B_x+B_y$ avec $B\le T$,
  donc $x$ est homogène. Ainsi $D_0,P',Q'$ sont homogènes, et $w$ est
  remplacé par sa composante de degré $k_W$.
\end{enumerate}
\end{document}
